\section{Coordenadas UV}

As coordenadas básicas são representadas por:

\begin{equation}
(u,v), \qquad 0\leq u\leq1,\quad 0\leq v\leq1.
\end{equation}

\(u\) representa a posição horizontal da textura e \(v\) representa a posição vertical. Os valores normalizados permitem acessar a textura independentemente de sua resolução.

\section{Translação}

A translação das coordenadas é dada por:

\begin{equation}
u'=u+\Delta u,\qquad v'=v+\Delta v.
\end{equation}

Os termos \(\Delta u\) e \(\Delta v\) determinam o deslocamento. Se o deslocamento for atualizado com o tempo, a textura pode aparentar movimento.

Uma forma simples é:

\begin{equation}
u'=u+s_ut,\qquad v'=v+s_vt,
\end{equation}

em que \(s_u\) e \(s_v\) representam velocidades de deslocamento.

\section{Escala}

A escala pode ser definida por:

\begin{equation}
u'=s_u u,\qquad v'=s_v v.
\end{equation}

Valores maiores que \(1\) aumentam a frequência aparente da textura, enquanto valores menores reduzem sua frequência.

\section{Rotação}

Considerando uma rotação em torno da origem do espaço UV:

\begin{equation}
u'=u\cos\theta-v\sin\theta,
\end{equation}

\begin{equation}
v'=u\sin\theta+v\cos\theta.
\end{equation}

O parâmetro \(\theta\) representa o ângulo de rotação.

\section{Transformação combinada}

As operações podem ser representadas matricialmente:

\begin{equation}
P'=M P,
\end{equation}

onde \(P=(u,v)^T\) é o ponto original e \(M\) representa a transformação escolhida.

\section{Interpolação}

Para um triângulo, as coordenadas UV dos fragmentos podem ser obtidas por interpolação baricêntrica:

\begin{equation}
UV=\alpha UV_0+\beta UV_1+\gamma UV_2,
\end{equation}

com:

\begin{equation}
\alpha+\beta+\gamma=1.
\end{equation}

Assim, os atributos dos vértices são utilizados para obter valores correspondentes em cada fragmento.

\section{Mapeamento procedural}

Uma formulação geral é:

\begin{equation}
UV'=f(P,t),
\end{equation}

ou, quando o objetivo é produzir diretamente a cor:

\begin{equation}
C=f(P,t).
\end{equation}

Essas expressões representam a ideia de que o resultado pode ser calculado a partir da posição espacial e, opcionalmente, do tempo.

\section{Relação com o pipeline}

O fluxo conceitual pode ser resumido como:

\begin{center}
\begin{tabular}{c}
\textit{Coordenadas UV}
$\rightarrow$
\textit{Transformação}
$\rightarrow$
\textit{Interpolação}
\\[0.2cm]
$\rightarrow$
\textit{Fragment Shader}
$\rightarrow$
\textit{Sampling}
$\rightarrow$
\textit{Cor final}
\end{tabular}
\end{center}